\appendix

\section{缩写与符号对照表}\label{sec:appendix}

表~\ref{tab:abbreviations}列出了本文中使用的主要缩写及其全称。

\begin{table}[htbp]
\centering
\caption{缩写与符号对照表}
\label{tab:abbreviations}
\begin{tabular}{ll}
\toprule
缩写 & 全称 \\
\midrule
MOF & Metal-Organic Framework（金属有机框架） \\
OWS & Overall Water Splitting（全分解水） \\
HER & Hydrogen Evolution Reaction（产氢反应） \\
OER & Oxygen Evolution Reaction（产氧反应） \\
AQE & Apparent Quantum Efficiency（表观量子效率） \\
AQY & Apparent Quantum Yield（表观量子产率） \\
PEC & Photoelectrochemical（光电化学） \\
LMCT & Ligand-to-Metal Charge Transfer（配体到金属电荷转移） \\
LSPR & Localized Surface Plasmon Resonance（局域表面等离子体共振） \\
DFT & Density Functional Theory（密度泛函理论） \\
BDC & 1,4-Benzenedicarboxylic Acid（对苯二甲酸） \\
TBAPy & 1,3,6,8-Tetrakis(4-benzoate)pyrene（四苯甲酸芘） \\
UiO & Universitetet i Oslo \\
MIL & Mat\'{e}riaux de l'Institut Lavoisier \\
ZIF & Zeolitic Imidazolate Framework（沸石咪唑酯框架） \\
HKUST & Hong Kong University of Science and Technology \\
MXene & 过渡金属碳化物/氮化物二维材料 \\
HOMO & Highest Occupied Molecular Orbital（最高占据分子轨道） \\
LUMO & Lowest Unoccupied Molecular Orbital（最低未占据分子轨道） \\
CB & Conduction Band（导带） \\
VB & Valence Band（价带） \\
\bottomrule
\end{tabular}
\end{table}

表~\ref{tab:symbols}列出了本文涉及的主要物理量符号。

\begin{table}[htbp]
\centering
\caption{物理量符号对照表}
\label{tab:symbols}
\begin{tabular}{lll}
\toprule
符号 & 物理量 & 单位 \\
\midrule
$\Delta G$ & 吉布斯自由能变 & $\mathrm{kJ \cdot mol^{-1}}$ \\
$E_g$ & 带隙能量 & $\mathrm{eV}$ \\
$\lambda$ & 光照波长 & $\mathrm{nm}$ \\
$r_{\mathrm{H_2}}$ & 产氢速率 & $\mathrm{mmol \cdot h^{-1} \cdot g^{-1}}$ \\
$S_{\mathrm{BET}}$ & BET比表面积 & $\mathrm{m^2 \cdot g^{-1}}$ \\
\bottomrule
\end{tabular}
\end{table}
